% GENERATED FILE. Edit canonical lessons, then run npm run generate:books.
% canonical-source-hash: fnv1a64:67f9b6fbede6df6a
% canonical-lessons: ZH-C234-tuichi, ZH-C234-quxiao, ZH-C234-gaitian, ZH-C234-yuyi, ZH-C234-duoyu

\chapter{Rain or Shine}
\label{ch:zh-rain-or-shine}
\hlchaptermodality{234}

\begin{chapteropening}
\textbf{By the end of this chapter:} \emph{I can put plans off, call them off, move them to another day, take a raincoat and shelter from the rain.}

Complete the last lesson of chapter 234: I can put plans off, call them off, move them to another day, take a raincoat and shelter from the rain.
\end{chapteropening}

\section[tuīchí]{tuīchí — to put off; to postpone}

\label{lesson:ZH-C234-tuichi}



\begin{quote}
Before the new one: say the Mandarin for slender, then the Mandarin for skin.
\end{quote}

\subsection*{You'll want to know: tuīchí}
\emph{tuīchí} — \textquotedblleft{}to put off; to postpone\textquotedblright{}.

\begin{quote}
\emph{míngtiān xiàyǔ wǒmen de yěcān yào tuīchí} — It's going to rain tomorrow, so our picnic will have to wait.
\end{quote}

Its characters come later; for now it is a word for the ear.

\begin{cousinweb}
\emph{tuīchí} is built from \emph{tuī} 'to push' and \emph{chí} 'slow', both inherited from Middle Chinese, \emph{thwoj} and \emph{drij} (Guangyun, in Baxter's transcription), which glosses the second 'slow' and 'leisurely'.
\end{cousinweb}

\begin{grammarlens}[title={What you've built}]
One more word for talking about plans.
\end{grammarlens}

\subsection*{Your turn}
\begin{itemize}
\raggedright
  \item \emph{Say it:} \emph{tuīchí}
  \item \emph{Say it:} \emph{tuīchí}, once more
  \item \emph{Say it:} \emph{míngtiān xiàyǔ wǒmen de yěcān yào tuīchí} — It's going to rain tomorrow, so our picnic will have to wait.
  \item \emph{Recall:} say the Mandarin for slender, then the Mandarin for skin, then say \emph{tuīchí} again
\end{itemize}

\begin{tcolorbox}[breakable,colback=teal!4,colframe=teal!35!black,title={Before you move on}]
What does tuīchí mean? (\textquotedblleft{}To put off; to postpone\textquotedblright{}.) Say it once more.
\end{tcolorbox}
\section[qǔxiāo]{qǔxiāo — to cancel; to call off}

\label{lesson:ZH-C234-quxiao}



\begin{quote}
Before the new one: say the Mandarin for skin, then the Mandarin for to put off.
\end{quote}

\subsection*{You'll want to know: qǔxiāo}
\emph{qǔxiāo} — \textquotedblleft{}to cancel; to call off\textquotedblright{}.

\begin{quote}
\emph{xiàyǔ le bǐsài qǔxiāo le} — It's raining, so the match is off.
\end{quote}

Its characters come later; for now it is a word for the ear.

\begin{cousinweb}
\emph{qǔxiāo} is built from \emph{qǔ} 'to take' and \emph{xiāo} 'to melt away', both inherited from Middle Chinese, \emph{tshjuX} and \emph{sjew} (Guangyun, in Baxter's transcription), which glosses the first 'to take in' and the second 'to wipe out' and 'to use up'.
\end{cousinweb}

\begin{grammarlens}[title={What you've built}]
One more word for talking about plans.
\end{grammarlens}

\subsection*{Your turn}
\begin{itemize}
\raggedright
  \item \emph{Say it:} \emph{qǔxiāo}
  \item \emph{Say it:} \emph{qǔxiāo}, once more
  \item \emph{Say it:} \emph{xiàyǔ le bǐsài qǔxiāo le} — It's raining, so the match is off.
  \item \emph{Recall:} say the Mandarin for skin, then the Mandarin for to put off, then say \emph{qǔxiāo} again
\end{itemize}

\begin{tcolorbox}[breakable,colback=teal!4,colframe=teal!35!black,title={Before you move on}]
What does qǔxiāo mean? (\textquotedblleft{}To cancel; to call off\textquotedblright{}.) Say it once more.
\end{tcolorbox}
\section[gǎitiān]{gǎitiān — some other day; another time}

\label{lesson:ZH-C234-gaitian}



\begin{quote}
Before the new one: say the Mandarin for to put off, then the Mandarin for to call off.
\end{quote}

\subsection*{You'll want to know: gǎitiān}
\emph{gǎitiān} — \textquotedblleft{}some other day; another time\textquotedblright{}.

\begin{quote}
\emph{jīntiān xiàyǔ wǒmen gǎitiān qù} — It's raining today; let's go another day.
\end{quote}

Its characters come later; for now it is a word for the ear.

\begin{cousinweb}
\emph{gǎitiān} is built from \emph{gǎi} 'to change' and \zh{天} 'day', both inherited from Middle Chinese, \emph{kojX} and \emph{then} (Guangyun, in Baxter's transcription), which glosses the first 'to alter'.
\end{cousinweb}

\begin{grammarlens}[title={What you've built}]
One more word for talking about plans.
\end{grammarlens}

\subsection*{Your turn}
\begin{itemize}
\raggedright
  \item \emph{Say it:} \emph{gǎitiān}
  \item \emph{Say it:} \emph{gǎitiān}, once more
  \item \emph{Say it:} \emph{jīntiān xiàyǔ wǒmen gǎitiān qù} — It's raining today; let's go another day.
  \item \emph{Recall:} say the Mandarin for to put off, then the Mandarin for to call off, then say \emph{gǎitiān} again
\end{itemize}

\begin{tcolorbox}[breakable,colback=teal!4,colframe=teal!35!black,title={Before you move on}]
What does gǎitiān mean? (\textquotedblleft{}Some other day; another time\textquotedblright{}.) Say it once more.
\end{tcolorbox}
\section[\texorpdfstring{\zh{雨衣} (yǔyī)}{雨衣 (yǔyī)}]{\zh{雨衣} (yǔyī) — raincoat}

\label{lesson:ZH-C234-yuyi}



\begin{quote}
Before the new one: say the Mandarin for to call off, then the Mandarin for some other day.
\end{quote}

\subsection*{You'll want to know: \zh{雨衣}}
\textbf{\zh{雨衣}} — \emph{yǔyī} — \textquotedblleft{}raincoat\textquotedblright{}.

\begin{quote}
\textbf{\zh{下雨了风也很大我有雨衣}} \emph{xiàyǔ le fēng yě hěn dà wǒ yǒu yǔyī} — It's raining and it's windy too; good thing I've got a raincoat.
\end{quote}

\begin{cousinweb}
\textbf{\zh{雨衣}} (\emph{yǔyī}) is built from \zh{雨} 'rain' and \zh{衣} 'garment', both inherited from Middle Chinese, \emph{hjuX} and \emph{'j+j} (Guangyun, in Baxter's transcription).
\end{cousinweb}

\begin{grammarlens}[title={What you've built}]
One more word for talking about plans.
\end{grammarlens}

\subsection*{Your turn}
\begin{itemize}
\raggedright
  \item \emph{Say it:} \emph{yǔyī}
  \item \emph{Say it:} \emph{yǔyī}, once more
  \item \emph{Say it:} \emph{xiàyǔ le fēng yě hěn dà wǒ yǒu yǔyī} — It's raining and it's windy too; good thing I've got a raincoat.
  \item \emph{Recall:} say the Mandarin for to call off, then the Mandarin for some other day, then say \emph{yǔyī} again
\end{itemize}

\begin{tcolorbox}[breakable,colback=teal!4,colframe=teal!35!black,title={Before you move on}]
What does \zh{雨衣} mean? (\textquotedblleft{}Raincoat\textquotedblright{}.) Say it once more.
\end{tcolorbox}
\section[duǒyǔ]{duǒyǔ — to shelter from the rain}

\label{lesson:ZH-C234-duoyu}



\begin{quote}
Before the new one: say the Mandarin for some other day, then the Mandarin for a raincoat.
\end{quote}

\subsection*{You'll want to know: duǒyǔ}
\emph{duǒyǔ} — \textquotedblleft{}to shelter from the rain\textquotedblright{}.

\begin{quote}
\emph{xiàyǔ le wǒmen qù fàndiàn duǒyǔ} — It started raining, so we sheltered in a restaurant.
\end{quote}

Its characters come later; for now it is a word for the ear.

Tone: duǒ and yǔ are both third tones, so the first rises and it sounds like duóyǔ.

\begin{cousinweb}
\emph{duǒyǔ} is built from \emph{duǒ} 'to hide' and \zh{雨} 'rain'. The second is inherited from Middle Chinese \emph{hjuX} (Guangyun, in Baxter's transcription); no dictionary we use records an older form of the first, so its history is not settled.
\end{cousinweb}

\begin{grammarlens}[title={What you've built}]
One more word for talking about plans.
\end{grammarlens}

\subsection*{Your turn}
\begin{itemize}
\raggedright
  \item \emph{Say it:} \emph{duǒyǔ}
  \item \emph{Say it:} \emph{duǒyǔ}, once more
  \item \emph{Say it:} \emph{xiàyǔ le wǒmen qù fàndiàn duǒyǔ} — It started raining, so we sheltered in a restaurant.
  \item \emph{Recall:} say the Mandarin for some other day, then the Mandarin for a raincoat, then say \emph{duǒyǔ} again
\end{itemize}

\begin{tcolorbox}[breakable,colback=teal!4,colframe=teal!35!black,title={Before you move on}]
What does duǒyǔ mean? (\textquotedblleft{}To shelter from the rain\textquotedblright{}.) Say it once more.
\end{tcolorbox}
